%==================================================================
% DATA TUGAS AKHIR — SATU-SATUNYA BERKAS YANG WAJIB DIISI
%
% Seluruh halaman sampul, persetujuan, pengesahan, dan pernyataan
% mengambil datanya dari sini. Tidak perlu mengetik ulang di tempat lain.
%==================================================================

%------------------------------------------------------------------
% BAHASA NASKAH
% id : istilah struktural Bahasa Indonesia (BAB, Gambar, Tabel)
% en : istilah struktural Bahasa Inggris (CHAPTER, Figure, Table)
%
% Halaman legal (persetujuan, pengesahan, pernyataan) tetap Bahasa
% Indonesia pada kedua pilihan, karena merupakan dokumen resmi institusi.
%------------------------------------------------------------------
\newcommand{\bahasanaskah}{id}

%------------------------------------------------------------------
% LABEL RUMUSAN MASALAH DAN TUJUAN
% Menentukan awalan label pada daftar rumusan masalah dan tujuan,
% sekaligus pada seluruh rujukan \cref di dalam naskah.
% Contoh: ubah menjadi P dan T, atau RM dan TP, sesuai kelaziman prodi.
%------------------------------------------------------------------
\newcommand{\labelrq}{}      % menghasilkan RQ1, RQ2, ...
\newcommand{\labeltujuan}{}   % menghasilkan T1, T2, ...

% Lebar kotak label. Bawaan 2.4em pas untuk "RQ1." dan "T1.".
% Bila kedua label di atas dikosongkan sehingga tampil sebagai nomor
% biasa, sempitkan menjadi 0.75em agar indentasinya sama dengan daftar
% bernomor lain di dalam naskah.
\newcommand{\lebarlabelrq}{2.4em}

%------------------------------------------------------------------
% GAYA SITASI
%
% apalike-id : APA nama-tahun dengan kata sambung "dan" (BAWAAN, teruji).
%              Berkas apalike-id.bst disertakan dalam template ini.
% apalike    : sama, tetapi kata sambungnya "and" (bawaan LaTeX).
% apacite    : APA yang lebih ketat. Perlu paket apacite, tersedia di
%              Overleaf dan MiKTeX penuh tetapi tidak pada semua
%              instalasi. Bila dipakai, aktifkan juga baris
%              \usepackage{apacite} pada xx-preambles.tex bagian 11.
%------------------------------------------------------------------
\newcommand{\gayasitasi}{apalike-id}
% \newcommand{\gayasitasi}{apalike}
% \newcommand{\gayasitasi}{apacite}

%------------------------------------------------------------------
% IDENTITAS
%------------------------------------------------------------------
\newcommand{\judulid}{Implementasi Diegetic Latency Masking pada NPC sebagai Mekanisme Penanganan Latensi Respons pada Game Eksplorasi Kampus}
\newcommand{\judulen}{Implementing Diegetic Latency Masking in NPCs as a Response Latency Mitigation Mechanism for a Campus Exploration Game}
\newcommand{\penulis}{Marco Benedictus Tirtaatmadja}
\newcommand{\nim}{202304560040}

\newcommand{\tipe}{Tugas Akhir}
\newcommand{\type}{Undergraduate Thesis}
\newcommand{\gelar}{Sarjana Komputer}
\newcommand{\prodi}{Sistem Informasi}
\newcommand{\fakultas}{Fakultas Biosains, Teknologi, dan Inovasi}
\newcommand{\universitas}{Universitas Katolik Indonesia Atma Jaya}
\newcommand{\gelaren}{Bachelor of Computer Science}
\newcommand{\prodien}{Information Systems Study Program}
\newcommand{\fakultasen}{School of Bioscience, Technology, and Innovation}
\newcommand{\universitasen}{Atma Jaya Catholic University of Indonesia}
\newcommand{\kota}{Jakarta}
\newcommand{\tahun}{2026}

%------------------------------------------------------------------
% DATA UNTUK FORMULIR RESMI (FR-UAJ)
% Hanya diperlukan bila formulir resmi ikut disertakan.
%------------------------------------------------------------------
\newcommand{\stid}{12023003387}              % St. ID mahasiswa
\newcommand{\kotaformulir}{Tangerang}        % kota pada formulir resmi
\newcommand{\tglujian}{3 Februari 2026}      % tanggal ujian/sidang TA
\newcommand{\periodeujian}{Februari / 2026}  % periode ujian (bulan / tahun)
\newcommand{\tglkoreksi}{Januari 2026}       % tanggal pada formulir pernyataan pembimbing

%------------------------------------------------------------------
% TANGGAL
%------------------------------------------------------------------
\newcommand{\tglpersetujuan}{15 Juni 2026}
\newcommand{\tglpersetujuandua}{15 Juni 2026}
\newcommand{\tglpengesahan}{15 Juli 2026}
\newcommand{\tglpernyataan}{15 Juli 2026}

%------------------------------------------------------------------
% DOSEN DAN PEJABAT
%------------------------------------------------------------------
\newcommand{\pembimbing}{Ir. Julius Victor Manuel Bata, S.Kom., M.T.}
\newcommand{\NIPpembimbing}{120181818}

% Jumlah pembimbing: 1 atau 2. Mengatur Lembar Persetujuan dan formulir FR-UAJ.
\newcommand{\jumlahpembimbing}{1}
\newcommand{\pembimbingdua}{Nama Pembimbing II, S.Kom., M.Kom.}
\newcommand{\NIPpembimbingdua}{000000000}

\newcommand{\ketuapenguji}{Nama Ketua Penguji, S.Kom., M.T.}
\newcommand{\sekretaris}{Nama Sekretaris Penguji, S.Kom., M.T.}
\newcommand{\penguji}{Nama Anggota Penguji, S.Kom., M.T.}

\newcommand{\koorprodi}{Gregorius Airlangga, S.Kom., M.T., Ph.D.}
\newcommand{\NIPkoorprodi}{000000000}
\newcommand{\dekan}{Nama Dekan, S.T., M.T., Ph.D.}
\newcommand{\NIPdekan}{000000000}

%------------------------------------------------------------------
% KATA KUNCI
%------------------------------------------------------------------
\newcommand{\katakunci}{diegetic latency masking, NPC behavior, response latency, campus exploration game.}
\newcommand{\keywords}{diegetic latency masking, perilaku NPC, latensi respons, game explorasi kampus.}
